\documentclass[11pt]{amsart}
\usepackage[margin=1in]{geometry}
\usepackage{amsmath,amssymb,amsthm}
\usepackage{longtable,booktabs,array}
\usepackage{url}
\usepackage[hidelinks]{hyperref}
\newtheorem{theorem}{Theorem}[section]
\newtheorem{proposition}[theorem]{Proposition}
\newtheorem{lemma}[theorem]{Lemma}
\newtheorem{corollary}[theorem]{Corollary}
\theoremstyle{definition}
\newtheorem{definition}[theorem]{Definition}
\theoremstyle{remark}
\newtheorem{remark}[theorem]{Remark}
\newcommand{\lean}[1]{\path{#1}}
\sloppy

\title[Factorization identities for the Fermat quartic fourfold]{Factorization Identities and Plane-Ideal Membership for the Fermat Quartic Fourfold, Formalized in Lean~4}
\author{Benjamin Stanley Frohman}
\email{frohmanbenjamin@gmail.com}
\date{November 4, 2026}
\begin{document}
\begin{abstract}
The Fermat quartic fourfold $X=\{z_0^4+\dots+z_5^4=0\}\subset\mathbb{P}^5$ contains the linear planes $\{z_0=\zeta z_1,\ z_2=\zeta z_3,\ z_4=\zeta z_5\}$ for every $\zeta$ with $\zeta^4=-1$. We formalize in Lean~4 with Mathlib, over an arbitrary commutative ring $R$, the polynomial identities behind this fact: the division of $z^4+w^4$ by $z-\zeta w$ with explicit remainder $(1+\zeta^4)w^4$, its vanishing when $\zeta^4=-1$, the four-factor identity $(x\pm\zeta y)(x\pm\zeta^3y)=x^4+y^4$, explicit three-term decompositions of the Fermat quartic, and membership of the Fermat quartic in the ideal generated by the three linear forms. We also prove the sextic analogue of the division identity, record the count $15\cdot4^3=960$, and state precisely which parts (spanning of Hodge classes, cycle classes) are \emph{not} formalized. All proofs are free of \texttt{sorry} and \texttt{native\_decide} and use only Lean's standard axioms.
\end{abstract}
\maketitle

\section{Introduction}
Linear subspaces on Fermat hypersurfaces are the classical source of algebraic cycles on these varieties; for the Fermat quartic fourfold, the planes obtained by pairing the six coordinates and choosing three fourth roots of $-1$ are known to span the primitive Hodge classes over $\mathbb{Q}$ (Aljovin--Movasati--Villaflor). The containment of each such plane in $X$ is an elementary algebraic identity. This note formalizes those identities exactly as stated below. Throughout, $R$ is an arbitrary commutative ring and all variables are elements of $R$; in particular the results hold for $R=k[z_0,\dots,z_5,\zeta]/(\zeta^4+1)$ and similar rings by instantiation, but no polynomial ring is constructed in the formalization.

\section{Division identities}
\begin{theorem}[\lean{FermatPlanes.pair_with_remainder}]\label{thm:pr}
For all $z,w,\zeta\in R$,
\[
z^4+w^4=(z-\zeta w)\,(z^3+\zeta z^2w+\zeta^2zw^2+\zeta^3w^3)+(1+\zeta^4)\,w^4.
\]
\end{theorem}
\begin{corollary}\label{cor:pv}
Let $\zeta^4=-1$. Then for all $z,w\in R$:
\begin{enumerate}
\item (\lean{FermatPlanes.pair_vanishing}) $z^4+w^4=(z-\zeta w)(z^3+\zeta z^2w+\zeta^2zw^2+\zeta^3w^3)$;
\item (\lean{FermatPlanes.pair_vanishing_neg}) $z^4+w^4=(z+\zeta w)(z^3-\zeta z^2w+\zeta^2zw^2-\zeta^3w^3)$.
\end{enumerate}
\end{corollary}

\begin{theorem}[Four-factor identity]\label{thm:ff}
For all $x,y,\zeta\in R$ (\lean{FermatPlanes.four_factor_with_remainder}),
\[
(x+\zeta y)(x-\zeta y)(x+\zeta^3y)(x-\zeta^3y)=x^4+y^4-(1+\zeta^4)\bigl(\zeta^2x^2y^2-(\zeta^4-1)y^4\bigr),
\]
and hence, if $\zeta^4=-1$ (\lean{FermatPlanes.four_factor_vanishing}),
$(x+\zeta y)(x-\zeta y)(x+\zeta^3y)(x-\zeta^3y)=x^4+y^4$.
\end{theorem}

\section{The Fermat quartic and the planes $Z_1$, $Z_2$}
Let $F=z_0^4+z_1^4+z_2^4+z_3^4+z_4^4+z_5^4$ (Lean: \lean{fermatPoly}) and, for $u,v\in R$, write $C_\zeta(u,v)=u^3+\zeta u^2v+\zeta^2uv^2+\zeta^3v^3$.
\begin{theorem}\label{thm:dec}
Assume $\zeta^4=-1$. For all $z_0,\dots,z_5\in R$:
\begin{enumerate}
\item (\lean{FermatPlanes.quartic_in_Z1_factors})
\[F=(z_0-\zeta z_1)C_\zeta(z_0,z_1)+(z_2-\zeta z_3)C_\zeta(z_2,z_3)+(z_4-\zeta z_5)C_\zeta(z_4,z_5);\]
\item (\lean{FermatPlanes.quartic_in_Z2_factors})
\[\begin{aligned}F={}&(z_0+\zeta z_1)(z_0^3-\zeta z_0^2z_1+\zeta^2z_0z_1^2-\zeta^3z_1^3)\\&+(z_2-\zeta z_3)C_\zeta(z_2,z_3)+(z_4-\zeta z_5)C_\zeta(z_4,z_5).\end{aligned}\]
\end{enumerate}
\end{theorem}

\begin{theorem}[Ideal membership; \lean{FermatPlanes.fermat_mem_Z1}]\label{thm:mem}
If $\zeta^4=-1$, then for all $z_0,\dots,z_5\in R$,
\[
F\in(z_0-\zeta z_1,\ z_2-\zeta z_3,\ z_4-\zeta z_5)\subseteq R .
\]
\end{theorem}
\begin{proof}
Immediate from Theorem~\ref{thm:dec}(1): each summand is a multiple of a generator.
\end{proof}
\begin{remark}
Taking $R$ to be a polynomial ring with $z_i$ the indeterminates and $\zeta$ a primitive eighth root of unity in the coefficient field, Theorem~\ref{thm:mem} says that the plane $\{z_0=\zeta z_1,z_2=\zeta z_3,z_4=\zeta z_5\}$ lies on $X$. This geometric reading is an instantiation and interpretation; the formal statement is exactly Theorem~\ref{thm:mem}.
\end{remark}

\section{The sextic analogue}
Let $c_6(z,w)=z^5+\zeta z^4w+\zeta^2z^3w^2+\zeta^3z^2w^3+\zeta^4zw^4+\zeta^5w^5$ (Lean: \lean{cofactor6}).
\begin{theorem}\label{thm:sex}
For all $z,w,\zeta\in R$, $z^6+w^6=(z-\zeta w)\,c_6(z,w)+(1+\zeta^6)w^6$ (\lean{FermatPlanes.pair6_with_remainder}); if $\zeta^6=-1$ then $z^6+w^6=(z-\zeta w)\,c_6(z,w)$ (\lean{FermatPlanes.pair6_vanishing}).
\end{theorem}

\section{Combinatorial labels}
The repository encodes a plane as a record consisting of a pairing of the six coordinates into three ordered pairs and three indices in $\mathrm{Fin}\,4$ naming roots of $\mu^4=-1$; $Z_1$ and $Z_2$ are the records with the standard pairing $(z_0,z_1),(z_2,z_3),(z_4,z_5)$ and root indices $(0,0,0)$, $(1,0,0)$.
\begin{proposition}\label{prop:lab}
\begin{enumerate}
\item (\lean{FermatPlanes.pairing_count_formula}) $15\cdot4^3=960$.
\item (\lean{FermatPlanes.Z1_is_standard}, \lean{FermatPlanes.Z2_is_standard}) $Z_1$ and $Z_2$ use the standard pairing.
\item (\lean{FermatPlanes.construct_id}, \lean{FermatPlanes.easyCl_is_section}, \lean{FermatPlanes.easyCl_eq_construct}) On $\mathbb{Q}\times\mathbb{Q}$, the map \lean{construct} is the identity, and \lean{easyCl}$(a,b)=(a,b)=$\lean{construct}$(a,b)$.
\end{enumerate}
\end{proposition}
Item~(1) is the arithmetic $15\cdot64=960$ ($15$ being the number of perfect matchings of six objects and $4^3$ the number of root choices); the formalization does not prove that these $960$ labels are the planes of $X$ or that they are distinct. Item~(3) is definitional: the ``cycle class'' \lean{easyCl} on the span of $[Z_1],[Z_2]$ is modelled as the identity map of $\mathbb{Q}^2$, and no cohomology is involved. The kernel also checks the auto-generated lemma \lean{FermatPlanes.Coord.ofNat_toCtorIdx} produced by \texttt{deriving} on the coordinate type.

\section{Scope}
Not formalized: projective geometry, the variety $X$ as a scheme, intersection numbers, Hodge classes, the spanning theorem of Aljovin--Movasati--Villaflor, or any statement about the Hodge conjecture.

\section{Correspondence between claims and Lean declarations}
All declarations are in namespace \lean{FermatPlanes}; Lean~4 v4.22.0 with Mathlib.
\begin{center}
\begin{tabular}{@{}p{0.25\textwidth}p{0.42\textwidth}p{0.28\textwidth}@{}}
\toprule
Claim & Lean declaration & File\\
\midrule
Theorem~\ref{thm:pr} & \lean{pair_with_remainder} & \lean{FermatIdentities.lean}\\
Corollary~\ref{cor:pv}(1) & \lean{pair_vanishing} & \lean{FermatIdentities.lean}\\
Corollary~\ref{cor:pv}(2) & \lean{pair_vanishing_neg} & \lean{FermatIdentities.lean}\\
Theorem~\ref{thm:ff} & \lean{four_factor_with_remainder}, \lean{four_factor_vanishing} & \lean{FourFactorIdentity.lean}\\
Theorem~\ref{thm:dec}(1) & \lean{quartic_in_Z1_factors} & \lean{FermatIdentities.lean}\\
Theorem~\ref{thm:dec}(2) & \lean{quartic_in_Z2_factors} & \lean{FermatIdentities.lean}\\
Theorem~\ref{thm:mem} & \lean{fermat_mem_Z1} & \lean{Membership.lean}\\
Theorem~\ref{thm:sex} & \lean{pair6_with_remainder}, \lean{pair6_vanishing} & \lean{SexticPair.lean}\\
Proposition~\ref{prop:lab}(1) & \lean{pairing_count_formula} & \lean{LinearPlanes.lean}\\
Proposition~\ref{prop:lab}(2) & \lean{Z1_is_standard}, \lean{Z2_is_standard} & \lean{EasyArrow.lean}\\
Proposition~\ref{prop:lab}(3) & \lean{construct_id}, \lean{easyCl_is_section}, \lean{easyCl_eq_construct} & \lean{FermatDatum.lean}, \lean{EasyArrow.lean}\\
\bottomrule
\end{tabular}
\end{center}

\appendix
\section{Verification certificate}\label{app:verification-certificate}

The formal development was machine-checked as follows.

{\raggedright
\begin{itemize}
\item Repository:\\ {\footnotesize\url{https://github.com/BenFrohman/FermatPlanes}}
\item Verified commit:\\ {\footnotesize\texttt{a8f4b5ec2ec53d7fcdf29bb37b79a961d4b460ba}}
\item Lean toolchain: \texttt{leanprover/lean4:v4.22.0}
\item Mathlib: tag \texttt{v4.22.0}, commit (from \texttt{lake-manifest.json})\\ {\footnotesize\texttt{79e94a093aff4a60fb1b1f92d9681e407124c2ca}}
\item Continuous integration run (build, axiom audit, \texttt{sorry}/\texttt{axiom} scan; passed):\\ {\footnotesize\url{https://github.com/BenFrohman/FermatPlanes/actions/runs/41794842261}}
\item SHA-256 of the GitHub source tarball of the verified commit, downloaded twice with identical result:\\ {\footnotesize\url{https://github.com/BenFrohman/FermatPlanes/archive/a8f4b5ec2ec53d7fcdf29bb37b79a961d4b460ba.tar.gz}}\\
{\footnotesize\texttt{cd2b240123e89b8e96939a3b53968fe5b44f10bdf974714c17096972fb1c6ed5}}
\end{itemize}\par}

The tarball hash refers to the verified commit rather than to a release tag, since a document cannot contain the hash of an archive that includes itself. The Lean sources are unchanged between that commit and the release tag.

\textbf{The development contains no \texttt{sorry} and no user-declared axioms.} Every one of the 17 theorems and lemmas listed below depends only on Lean's standard axioms \texttt{propext}, \texttt{Classical.choice} and \texttt{Quot.sound} (or on none). The table is the verbatim output of \texttt{lake env lean AxiomAudit.lean} (\texttt{\#print axioms}) on the verified commit.

{\footnotesize
\begin{longtable}{p{0.55\linewidth}p{0.38\linewidth}}
\hline
Lean declaration & \texttt{\#print axioms} \\
\hline\endhead
\path{FermatPlanes.Z1_is_standard} & \texttt{[propext]} \\
\path{FermatPlanes.Z2_is_standard} & \texttt{[propext]} \\
\path{FermatPlanes.easyCl_eq_construct} & \texttt{[propext]} \\
\path{FermatPlanes.easyCl_is_section} & \texttt{[propext]} \\
\path{FermatPlanes.construct_id} & \texttt{[propext]} \\
\path{FermatPlanes.pair_vanishing} & \texttt{[propext, Quot.sound]} \\
\path{FermatPlanes.pair_vanishing_neg} & \texttt{[propext, Quot.sound]} \\
\path{FermatPlanes.pair_with_remainder} & \texttt{[propext, Quot.sound]} \\
\path{FermatPlanes.quartic_in_Z1_factors} & \texttt{[propext, Quot.sound]} \\
\path{FermatPlanes.quartic_in_Z2_factors} & \texttt{[propext, Quot.sound]} \\
\path{FermatPlanes.four_factor_vanishing} & \texttt{[propext, Quot.sound]} \\
\path{FermatPlanes.four_factor_with_remainder} & \texttt{[propext, Quot.sound]} \\
\path{FermatPlanes.Coord.ofNat_toCtorIdx} & \texttt{(no axioms)} \\
\path{FermatPlanes.pairing_count_formula} & \texttt{(no axioms)} \\
\path{FermatPlanes.fermat_mem_Z1} & \texttt{[propext, Quot.sound]} \\
\path{FermatPlanes.pair6_vanishing} & \texttt{[propext, Quot.sound]} \\
\path{FermatPlanes.pair6_with_remainder} & \texttt{[propext, Quot.sound]} \\
\hline
\end{longtable}}

\begin{thebibliography}{9}
\bibitem{amv} J.~Aljovin, H.~Movasati, R.~Villaflor~Loyola, \emph{Integral Hodge conjecture for Fermat varieties}, J.\ Symbolic Comput.\ \textbf{95} (2019), 173--194.
\bibitem{mathlib} The mathlib Community, \emph{The Lean mathematical library}, CPP 2020, 367--381.
\end{thebibliography}
\end{document}
